% !TEX root = ../Projektdokumentation.tex
\section{Fazit} 
\label{sec:Fazit}
In diesem Abschnitt geht es darum den Abschluss des Projektes sachlich zu bewerten und transparent darzustellen,
inwiefern die Projektplanung eingehalten werden konnte. Außerdem wird gezeigt, was der Autor gelernt hat und
wie sich das im Projekt behandelte System in Zukunft weiterentwickeln wird.


\subsection{Soll-/Ist-Vergleich}
\label{sec:SollIstVergleich}
Das Projektziel die dezentralisierte Verfügbarkeit von Daten zweier Kundendaten speichernder Systeme aufzulösen und 
somit eine Source of Truth zu schaffen wurde erreicht.
Die Funktionalität und Einhaltung der Anforderungen wurde somit erfüllt. 
\par\smallskip
Bei der Einhaltung der Projektplanung ist es im Hinblick auf Zeit und Kosten zu Abweichungen gekommen. 
So wurde das Zeitlimit, um fünf Stunden überschritten und die Kosten sind dazu proportional um 175 Euro gestiegen. 
Grund ist, dass die Umstellung vom im Projektantrag erwähnten Hangfire-System auf ein EventQueue-basiertes System 
wie RabbitMQ neue Herausforderungen mit sich brachte. Zum einen war die Technologie im Unternehmen bisher nicht im Einsatz, 
wodurch ein erhebliches mehr in Einarbeitung geleistet werden musste.
Zum anderen wurde klar, dass die Implementierung in diesem Prototypen als Beispielimplementierung für andere unternehmensinterne Projekte dienen kann,
weswegen das Projekt mehr Review-Zyklen als üblich unterstand. 
\par\smallskip
Wie in Tabelle~\ref{tab:Vergleich} zu erkennen ist, konnte die Zeitplanung nicht eingehalten werden.
\tabelle{Soll-/Ist-Vergleich}{tab:Vergleich}{Zeitnachher.tex}


\subsection{Lessons Learned}
\label{sec:LessonsLearned}
Der Prüfling hat bei der Durchführung des Projektes gelernt realistischere Einschätzungen zu machen in Bezug
auf die Zeitplanung bei bisher unbekannten Technologien.
Außerdem ist dem Prüfling nun der Umgang mit RabbitMQ, den einhergehenden Topologien und Einstellungen 
für Queues vertraut.
Zudem kann der Prüfling zukünftlich für Fragen in diesem Bereich zur Seite stehen und anderen Entwicklern bei der 
Implementierung helfen.  


\subsection{Ausblick}
\label{sec:Ausblick}
Im Rahmen des Projektes wurde die Anzahl der zu synchronisierenden Datenfelder bewusst begrenzt.
Zukünftig kann die Anwendung erweitert werden, sodass weitere kundenspezifische Daten zwischen den beteiligten Systemen synchronisiert werden.
Dadurch könnte der manuelle Pflegeaufwand in weiteren Prozessen zusätzlich reduziert werden.
\par\smallskip
Eine weitere mögliche Erweiterung besteht darin, die gespeicherte Synchronisationshistorie gezielter auszuwerten.
Über eine interne Schnittstelle oder ein Dashboard könnten beispielsweise fehlgeschlagene Synchronisationen,
wiederholte Verarbeitungsversuche oder häufig auftretende Konflikte sichtbar gemacht werden.
Dadurch könnten technische Fehler schneller erkannt und die Datenqualität langfristig besser überwacht werden.